\documentclass[../proyecto.tex]{subfiles}
\graphicspath{{\subfix{../images/}}}
\begin{document}

\section{Metodología}

% máximo 2 páginas

\subsection{Análisis de referentes}

Se realizará una revisión bibliográfica sobre referentes artísticos, académicos, comerciales, de hardware y software libre.

Esta revisión se enfocará en temas relacionados a la síntesis de sintetizadores de forma popular, que permitirá identificar conceptos y debates revelantes que permitirán contextualizar la investigación y práctica artística realizada.

\subsection{Prototipado electrónico}

La práctica electrónica se hará de forma iterativa, creando progresivamente distintos prototipos de sintetizadores y sus módulos constituyentes. Su socialización y su uso permitirá el aprendizaje a partir de errores, la escritura de estándares y estrategias de fabricación y uso, que son una práctica común en la comunidad de hardware libre, y que permitirá la creación de sintetizadores con distintas variantes.

\subsection{Prototipado computacional}

La práctica computacional se hará de forma rumiante, publicando productos de mínima viabilidad, y compartiendo el proceso de desarrollo en repositorios públicos, con el fin de que puedan todos los artefactos de la investigación ser analizados, replicados, adaptados y mejorados por pares artistas, académicas, y audiencias interesadas, a través del tiempo y de forma decentralizada.

\subsection{Escritura de documentación}

\subsection{Socialización de la investigación}

Durante la investigación doctoral se desarrollarán distintas herramientas que faciliten la síntesis de sintetizadores, los que serán prototipados, documentados y puestos a disposición de pares artistas, y serán usados en contextos de producción de obra artística, como material docente para la enseñanza de electrónica y computación, y para su posterior comercialización para su difusión y financiamiento de la investigación material realizada.

\subsection{Escritura y enseñanza de material base}

Para poder abordar de forma sostenida una investigación doctoral profunda en la síntesis de sintetizadores, se están desarrollando en paralelo cursos de pregrado para diseño que permitan a estudiantes de pregrado poner en práctica lo investigado, y que esto a su vez nutra de nuevas ideas y experiencias a esta investigación. Como subproducto de esta enseñanza se están escribiendo libros de referencia de enseñanza de electrónica y computación aplicadas a sonido, que permitan un diálogo constante con audiencias de estudiantes, artistas e investigadoras.

\end{document}

